\section{Introduction}

Denoising-based generative models, including diffusion models~\citep{ho2020ddpm,song2021score,song2020sde} and flow matching models~\citep{lipman2022flow}, have achieved remarkable success in high-fidelity image generation. Among them, Diffusion Transformers (DiTs) have become a particularly scalable backbone due to their strong modeling capacity and compatibility with large-scale training. Despite these advances, training DiTs remains computationally expensive, and improving their optimization efficiency without sacrificing generation quality is still an important challenge.

A growing body of work suggests that the intermediate representations of diffusion transformers play a critical role in both training dynamics and final generation quality. Representation alignment methods such as REPA~\citep{yu2024repa} show that aligning DiT hidden states with strong visual representations from pretrained encoders can substantially accelerate training. Subsequent studies further indicate that the benefit of such alignment is closely tied to the structure and quality of intermediate features, motivating a broader view that improving internal representations can be an effective route to improving generative modeling. However, these methods often rely on external representation encoders, which introduce additional computation, architectural dependency, and reduced portability to domains where strong pretrained teachers are unavailable.

Recent self-contained methods attempt to remove this dependency by using the diffusion transformer itself as a source of representation guidance. SRA~\citep{jiang2025sra} observes that representations in diffusion transformers tend to improve as noise decreases and depth increases, and therefore aligns earlier or noisier features with later or cleaner features from an EMA teacher. LayerSync~\citep{haghighi2026layersync} similarly exploits cross-layer structure by encouraging shallow representations to synchronize with deeper, more informative ones. These works suggest that useful supervision already exists inside the model: different layers and timesteps provide different views of the same denoising process, and this internal structure can guide representation learning without an external teacher.

Nevertheless, directly aligning hidden states is only one way to exploit cross-layer information. Inspired by shortcut models and consistency-style learning, we ask whether the evolution of DiT representations across depth can itself be used as a training signal. Rather than simply forcing two layers to be close, we learn a conditional transition between them: given a hidden state at an earlier layer, a lightweight predictor estimates the corresponding representation at a later layer. This turns the DiT depth trajectory into a self-supervised representation learning problem. A useful intermediate representation should not only support the final denoising or velocity prediction, but should also evolve in a predictable and compositional manner across transformer depth.

To this end, we propose \emph{Depth Shortcut Prediction} (DSP), a self-contained auxiliary objective for diffusion transformers. DSP trains a lightweight predictor to map a source-layer hidden state to a target-layer hidden state conditioned on the source layer, target layer, and timestep. To handle the scale variation of hidden states, we formulate prediction in direction--magnitude space: the predictor estimates both the target token directions and the residual change in log magnitude. The training objective consists of two auxiliary losses, one for direction prediction and one for magnitude prediction. Each loss combines direct cross-depth supervision with a bootstrap consistency term, encouraging composed short transitions to agree with direct long transitions.

Our method is designed primarily as a representation learning objective. By supervising cross-depth prediction over many layer pairs, DSP encourages the DiT backbone to learn intermediate features whose evolution is structured, predictable, and useful for generation. As a secondary consequence, the learned predictor can also reconstruct skipped hidden states and enable optional layer-skipping inference, providing a controllable quality--speed tradeoff. However, unlike methods focused purely on efficient inference, our main goal is to improve representation quality and generative training using the model's own hidden-state trajectory.

Our main contributions are as follows:
\begin{itemize}
    \item We introduce Depth Shortcut Prediction, a self-supervised representation learning objective that uses cross-depth hidden-state prediction as internal guidance for diffusion transformers.

    \item We formulate hidden-state prediction in direction--magnitude space, allowing the model to learn both target representation orientation and scale through two auxiliary losses.

    \item We incorporate bootstrap depth consistency, encouraging learned transitions to be compositional across transformer depth rather than only fitting isolated layer pairs.

    \item We show that the resulting objective improves DiT representation learning and generation quality, while also enabling optional layer-skipping inference as an auxiliary benefit.
\end{itemize}